%% kapa: exam=1 round=21 year=2010 period=2 form=A subject=appraisal_law subject_ko=부동산관계법규 seq=5 number=45
%% source: Q-Net 감정평가사 기출문제 (cst00309) artlSeq=5024181 file=extracted/1008/5024181/2010년 제21회 감정평가사 1차시험 기출문제/2교시.hwp
%% conversion: kordoc 4.24.51 -> tools/gen1.py (auto); review=proofread-gpt-sol; flags=none
\begin{question}{45}
국토의 계획 및 이용에 관한 법령상 토지거래허가에 관한 설명으로 옳은 것은? (단, 다툼이 있으면 판례에 의함)
\begin{choices}
  \item 국토해양부장관은 토지거래허가구역의 지정기간이 끝나는 허가구역을 계속하여 다시 허가구역으로 지정할 때 그 지정이 시급한 경우에는 중앙도시계획위원회의 심의를 거칠 필요가 없다.
  \item 토지거래허가구역의 지정에 대하여 항고소송을 제기할 수 없다.
  \item 토지거래허가구역 내 토지에 관하여 매매계약을 체결하고 계약금만 주고받은 상태에서 토지거래계약허가를 받은 경우, 매도인은 민법규정에 따라 그 계약을 해제할 수 있다.
  \item 토지거래허가구역 지정기간 중에 허가구역 안의 토지에 대하여 토지거래계약허가를 받지 않고 토지거래계약을 체결한 후 허가구역 지정이 해제된 경우에도, 토지거래계약허가를 받아야 한다.
  \item 토지거래허가구역에 거주하는 농업인이 그 허가구역 내에서 농업을 경영하기 위하여 토지거래계약허가를 받은 경우에는 토지의 취득시부터 5년 동안 그 토지를 허가받은 목적대로 이용하여야 한다.
\end{choices}
\end{question}
\answer{3}
